%% Revised working manuscript for Results in Applied Mathematics.
%% Historical sources and results are preserved outside this revision.
\documentclass[a4paper,fleqn]{cas-sc}
\usepackage[super,sort&compress]{natbib}
\usepackage{amsmath,amssymb,amsthm,mathtools,booktabs,tabularx}
\usepackage{algorithm,algpseudocode}
\usepackage{hyperref}
\hypersetup{hypertexnames=false}
% CAS 2.4 prints an empty ORCID footnote when identifiers are not supplied.
\renewcommand{\printorcid}{}
% Preserve the keyword layout with a deliberate, non-overfull overlay.
\ExplSyntaxOn
\RenewDocumentEnvironment{Abstract}{o}
{
  \group_begin:
  \IfNoValueTF{#1}{}{\tex_gdef:D\abstractname{#1}}
  \parindent0pt
  \box_if_empty:NTF\g_stm_key_box
    {\leftskip=.35\textwidth}
    {
      \dim_set:Nn\l_tmpa_dim{\box_ht:N\g_stm_key_box+\box_dp:N\g_stm_key_box}
      \leftskip=.35\textwidth
      \hspace*{-.35\textwidth}
      \noindent\hbox_to_wd:nn{0pt}{\box\g_stm_key_box\hss}
      \skip_vertical:n{-\l_tmpa_dim}
    }
  \noindent\abstractname\par
  \skip_vertical:n{-4pt}
  \noindent\rule{.65\textwidth}{.2pt}\par\footnotesize
  \ignorespaces\everypar{\parindent=1.5em}
}{\par\group_end:}
\ExplSyntaxOff
\newcommand{\R}{\mathbb R}
\newcommand{\norm}[1]{\left\lVert#1\right\rVert}
\newcommand{\ip}[2]{\left\langle#1,#2\right\rangle}
\DeclareMathOperator{\rank}{rank}
\DeclareMathOperator{\dist}{dist}
\DeclareMathOperator{\diag}{diag}
\DeclareMathOperator*{\argmin}{arg\,min}
\newtheorem{theorem}{Theorem}
\newtheorem{proposition}[theorem]{Proposition}
\newtheorem{lemma}[theorem]{Lemma}
\newtheorem{corollary}[theorem]{Corollary}
\theoremstyle{remark}
\newtheorem{remark}{Remark}
\begin{document}
\let\WriteBookmarks\relax
\shorttitle{Graph energy bounds and nuclear-norm budgets}
\shortauthors{C. Ram\'irez Ovalle et al.}
\title[mode=title]{Graph Energy Bounds and Global Nuclear-Norm Budgets in Fully Connected Neural Networks}
\author[1]{Carlos Ram\'irez Ovalle}
\cormark[1]
\ead{carlosovalle@javerianacali.edu.co}
\author[1]{Abel Alvarez}
\ead{abel.alvarez@javerianacali.edu.co}
\author[1]{Andr\'es J. Castrill\'on V\'asquez}
\ead{ajcastrillon@javerianacali.edu.co}
\affiliation[1]{organization={Pontificia Universidad Javeriana Cali},city={Cali},country={Colombia}}
\cortext[1]{Corresponding author}
\begin{abstract}
The spectral energy of a layered neural-network graph and an additive nuclear-norm budget describe related but distinct quantities. For signed weights connecting consecutive layers, we prove that the nuclear norm of the full symmetric adjacency matrix lies between one half and one times the sum of the individual bipartite-block energies. A perturbation estimate yields a local square-root loss-gap bound under quadratic growth relative to a closed minimizer set, with strong convexity as a special case. We derive the Euclidean projection onto a global nuclear-norm budget through a common threshold on all layer singular values, and distinguish it from separate layer budgets, spectral-norm projection and spectral normalization. For a split constrained formulation, we establish finite-iteration feasibility bounds and, under explicit smoothness and projection assumptions, an a posteriori stationarity bound. These certificates do not imply convergence of stochastic nonconvex ADMM. Counterexamples delimit stronger claims about penalty gaps and spectral control. The mathematical statements are supported by small reproducible numerical checks. Historical neural-network experiments are retained as unaudited exploratory evidence; a controlled evaluation on MNIST and Fashion-MNIST, including projected stochastic gradient and matched-budget baselines, remains to be completed. The contribution is a precise connection between layered graph energy, parameter-space stability and verifiable spectral-budget diagnostics, rather than a new nuclear-norm optimization primitive or a demonstrated predictive advantage.
\end{abstract}
\begin{keywords}
Graph energy \sep Nuclear norm \sep Quadratic growth \sep Euclidean projection \sep ADMM \sep Neural networks
\end{keywords}
\maketitle
\noindent\textbf{Revision status.} Mathematical working revision for \emph{Results in Applied Mathematics}, 13 October 2026. The replacement neural-network evaluation and author approval of the declarations are pending. This document is not a submission-ready empirical report.
\section{Introduction}\label{sec:intro}
A dense feedforward network carries a natural signed weighted graph: neurons are vertices and matrix entries are edge weights. The sum of the absolute eigenvalues of a symmetric adjacency matrix is its graph energy, equivalently its nuclear norm. The classical graph-energy literature provides the spectral starting point~\cite{Gutman2012}. Extending this definition to real signed weights is convenient for neural networks, but does not identify spectral energy with electrical consumption, inference operations, storage, or the complexity of the represented function.

Nuclear-norm regularization and singular-value thresholding are established tools~\cite{Watson1992Subdiff,Cai2010SVT}. Compression-aware training using low-rank regularization was studied by Alvarez and Salzmann~\cite{AlvarezSalzmann2017}; norm-constrained training using Frank--Wolfe was studied by Zimmer, Spiegel and Pokutta~\cite{Zimmer2022}. Consequently, neither nuclear-norm training nor imposing norm constraints on neural networks is new. The present contribution is an explicit derivation of graph-energy comparison bounds for consecutive-layer architectures, a consistent local stability statement, and finite-iteration diagnostics for a global spectral budget. We make no priority claim for the underlying matrix inequalities, projection method, or ADMM splitting.

The distinction between a graph quantity and its additive surrogate matters both mathematically and computationally. A single bipartite block has energy twice the nuclear norm of its weight matrix; overlapping blocks in a multilayer adjacency matrix generally do not have additive energies. The feasible set considered here budgets the additive surrogate. Its exact Euclidean projection couples the layers through a single singular-value threshold. Fixed layer budgets instead specify a usually smaller feasible set, which can be useful but solves a different projection problem.

Local quadratic growth bounds distance to a minimizer set by the square root of the loss gap. This implication provides conditional stability near that set; it does not show that a low spectral budget guarantees low loss, or that a trained ReLU network satisfies quadratic growth. Relations among quadratic growth, strong convexity and the Polyak--\L{}ojasiewicz condition require their own hypotheses~\cite{Karimi2016QGPL}. Our proofs use quadratic growth directly and do not infer it from selected positive Hessian eigenvalues.

The optimization analysis likewise separates an exact mathematical step from stochastic implementation. Standard ADMM is a useful splitting framework~\cite{Boyd2011ADMM}, but its convex convergence results do not automatically apply to inexact neural-network training with momentum, changing penalties and varying budgets. We therefore state finite-iteration feasibility and conditional stationarity certificates, and specify the experiments needed to assess practical behavior without assuming that ADMM is superior.

\section{Layered graph energy and its additive surrogate}\label{sec:graph}
Let $m\ge1$ be the number of weight matrices and let the $m+1$ neuron layers have positive widths $n_0,\ldots,n_m$. We use column activations and
\[
 W_\ell\in\R^{n_\ell\times n_{\ell-1}},\qquad
 b_\ell\in\R^{n_\ell},\qquad
 x_\ell=\phi_\ell(W_\ell x_{\ell-1}+b_\ell).
\]
The graph uses signed entries of $W_\ell$ without taking absolute values. Biases are trainable parameters but are excluded from the graph and budget. Adding constant vertices to encode biases would define another graph and another budget; we do not make that substitution implicitly.

The symmetric adjacency $A(W)$ has zero diagonal blocks and
\begin{equation}\label{eq:adjacency}
 A(W)_{\ell,\ell-1}=W_\ell,\quad
 A(W)_{\ell-1,\ell}=W_\ell^\top\quad(\ell=1,\ldots,m),
\end{equation}
with all other blocks zero. Even-indexed and odd-indexed neuron layers form a bipartition of the full graph. Define
\begin{equation}\label{eq:omega}
 E(A(W))=\norm{A(W)}_*,\qquad
 \Omega(W)=2\sum_{\ell=1}^m\norm{W_\ell}_*.
\end{equation}
Here $\norm{\cdot}_*$, $\norm{\cdot}_2$ and $\norm{\cdot}_F$ are respectively the nuclear, spectral and Frobenius matrix norms.

\begin{proposition}[Block identity and two-sided comparison]\label{prop:graph}
For the consecutive-layer adjacency~\eqref{eq:adjacency},
\begin{equation}\label{eq:bilateral}
 E\!\left(\begin{bmatrix}0&W_\ell^\top\\ W_\ell&0\end{bmatrix}\right)
 =2\norm{W_\ell}_*,\qquad
 \frac{\Omega(W)}2\le E(A(W))\le\Omega(W).
\end{equation}
These statements include zero weights and arbitrary real signs. No claim is made that the lower comparison constant is optimal.
\end{proposition}
\begin{proof}
If $W_\ell=Q_\ell\Sigma_\ell V_\ell^\top$ is a compact SVD, each positive singular value $\sigma_i$ yields eigenvectors $(v_i,\pm q_i)/\sqrt2$ of the displayed block with eigenvalues $\pm\sigma_i$. The other eigenvalues are zero. This proves the block identity. Write $A=\sum_\ell\widetilde B_\ell$, where each $\widetilde B_\ell$ is the corresponding block embedded in the full vertex space. The nuclear-norm triangle inequality gives the upper bound.

For the lower bound, put $\varepsilon_\ell=(-1)^{\lfloor\ell/2\rfloor}$ for layer indices $\ell=0,\ldots,m$ and $J=\diag(\varepsilon_0 I_{n_0},\ldots,\varepsilon_m I_{n_m})$. The product $\varepsilon_{\ell-1}\varepsilon_\ell$ is $+1$ for odd $\ell$ and $-1$ for even $\ell$. Thus
\[
 A_o=(A+JAJ)/2=\sum_{\ell\ \text{odd}}\widetilde B_\ell,
 \qquad A_e=(A-JAJ)/2=\sum_{\ell\ \text{even}}\widetilde B_\ell.
\]
Blocks of the same parity have disjoint vertex supports. Up to a simultaneous row and column permutation, each matrix is a direct sum of its bipartite blocks and zeros. Hence $\norm{A_o}_*+\norm{A_e}_*=\Omega(W)$. Orthogonal invariance and the triangle inequality give $\norm{A_o}_*\le\norm A_*$ and $\norm{A_e}_*\le\norm A_*$, yielding $\Omega(W)\le2\norm A_*$. Empty parity classes contribute zero.
\end{proof}

For two scalar weights $W_1=W_2=[1]$, the three-vertex path has eigenvalues $0,\pm\sqrt2$, so $E(A)=2\sqrt2$ and $\Omega=4$. The block identity is therefore not a full-network identity. The comparison proof depends on consecutive-layer connections; additional skip, recurrent or within-layer edges require a separate analysis.

\paragraph{Dependence on parametrization.}
For a two-layer bias-free ReLU network $h(x)=W_2\operatorname{ReLU}(W_1x)$ and any $c>0$, positive homogeneity gives
\[
 (c^{-1}W_2)\operatorname{ReLU}(cW_1x)=h(x),\qquad
 \Omega(cW_1,c^{-1}W_2)=2(c\norm{W_1}_*+c^{-1}\norm{W_2}_*).
\]
The latter generally changes with $c$; for scalar unit weights it changes from $4$ to $5$ at $c=2$. Predictions and data loss are unchanged, but weight decay generally changes. Thus $\Omega$ describes parameters, not an invariant of the represented function.

\section{Local stability relative to a minimizer set}\label{sec:stability}
The parameter space is the finite-dimensional Euclidean product
\[
 \Theta=\prod_{\ell=1}^m\R^{n_\ell\times n_{\ell-1}}
 \times\prod_{\ell=1}^m\R^{n_\ell},\quad
 \norm{\theta}_2^2=\sum_\ell\norm{W_\ell}_F^2+\sum_\ell\norm{b_\ell}_2^2.
\]
Write $\norm W_F^2=\sum_\ell\norm{W_\ell}_F^2$ and
$d_\ell=\min(n_\ell,n_{\ell-1})$, $D=\sum_\ell d_\ell$.

\begin{lemma}[Perturbation estimate]\label{lem:perturbation}
For any two weight collections $W,V$, let $R(W,V)=\sum_\ell\rank(W_\ell-V_\ell)$. Then
\begin{align}\label{eq:perturbation}
 |\Omega(W)-\Omega(V)|
 &\le 2\sum_\ell\norm{W_\ell-V_\ell}_*
 \le2\sqrt{R(W,V)}\norm{W-V}_F\\
 &\le2\sqrt D\norm{W-V}_F.\nonumber
\end{align}
For $\theta=(W,b)$ and $\eta=(V,c)$, $\norm{W-V}_F\le\norm{\theta-\eta}_2$.
\end{lemma}
\begin{proof}
The scalar triangle inequality and reverse triangle inequality for the nuclear norm give the first bound. For a matrix $M$ of rank $r$, Cauchy--Schwarz on its nonzero singular values gives $\norm M_*\le\sqrt r\norm M_F$. Applying this to each difference and then applying Cauchy--Schwarz across layers gives the second bound. Finally $R(W,V)\le D$ and the parameter norm includes the nonnegative squared bias differences. The rank-zero case is immediate.
\end{proof}
The pair-dependent $R(W,V)$ is not a uniform constant. An empirical numerical rank may underestimate the mathematical rank, so a thresholded rank cannot replace $R$ in a rigorous certificate. Known analytic ranks or the conservative dimension bound $D$ are suitable. The same proof, using $|\norm A_*-\norm{\widetilde A}_*|\le\norm{A-\widetilde A}_*$, also bounds $|E(A(W))-E(A(V))|$ by the right-hand sides of~\eqref{eq:perturbation}.

\begin{lemma}[Existence and projection]\label{lem:existence}
If $f_{\rm data}:\Theta\to\R$ is continuous and bounded below, and $\beta>0$, then
$f_{\rm reg}(\theta)=f_{\rm data}(\theta)+(\beta/2)\norm\theta_2^2$
has a nonempty compact global minimizer set. Every nonempty closed set $S\subset\Theta$ admits a nearest point to each $\theta\in\Theta$. That nearest point is unique when $S$ is convex; uniqueness is not asserted for a general minimizer set.
\end{lemma}
\begin{proof}
If $f_{\rm data}\ge a$, then $f_{\rm reg}(\theta)\ge a+(\beta/2)\norm\theta_2^2$, proving coercivity. A nonempty sublevel set is compact, and continuity yields an attained minimum and a compact minimizer set. For a closed, possibly unbounded $S$, fix $s_0\in S$. A minimizing sequence for distance from $\theta$ can be restricted to the closed ball of radius $\norm{\theta-s_0}_2+1$ about $\theta$. Compactness provides a convergent subsequence with limit in $S$ attaining the distance. Strict convexity of squared distance gives uniqueness on a convex set.
\end{proof}
Continuity alone is insufficient for the first conclusion: $f_{\rm data}(t)=-t^4$ defeats quadratic regularization. Regularization must cover all parameters whose escape is to be excluded. A weight budget bounds $W$ but does not make $(W,b)$ compact when biases are free. Existence does not establish quadratic growth.

\begin{theorem}[Quadratic-growth stability]\label{thm:qg}
Fix an objective $f:\Theta\to\R$ with finite attained global minimum $f_*$ and nonempty closed global minimizer set $S$. Suppose a specified neighborhood $\mathcal U$ of $S$ and $\gamma>0$ satisfy
\begin{equation}\label{eq:qg}
 f(\theta)-f_*\ge\frac\gamma2\dist(\theta,S)^2,
 \qquad\theta\in\mathcal U.
\end{equation}
For every $\theta=(W,b)\in\mathcal U$ and every nearest point $\theta^+=(W^+,b^+)\in S$,
\begin{equation}\label{eq:stability}
 |\Omega(W)-\Omega(W^+)|
 \le2\sqrt{\frac{2R(W,W^+)}\gamma}\sqrt{f(\theta)-f_*}
 \le2\sqrt{\frac{2D}\gamma}\sqrt{f(\theta)-f_*}.
\end{equation}
The same inequalities hold with the left side replaced by $|E(A(W))-E(A(W^+))|$.
\end{theorem}
\begin{proof}
Lemma~\ref{lem:existence} supplies a nearest point. By Lemma~\ref{lem:perturbation},
\[
 |\Omega(W)-\Omega(W^+)|\le2\sqrt{R(W,W^+)}\norm{W-W^+}_F
 \le2\sqrt{R(W,W^+)}\dist(\theta,S).
\]
Apply~\eqref{eq:qg} and $R\le D$. The adjacency version uses the observation after Lemma~\ref{lem:perturbation}.
\end{proof}

The objective in this theorem is fixed once and for all: choosing $f=f_{\rm reg}$ does not provide a certificate for a gap in $f_{\rm data}$ or in a hinge-penalized objective. A small loss gap implies a bounded spectral variation under~\eqref{eq:qg}; the reverse implication has not been established. Nor is $\Omega(W^+)$ necessarily the smallest surrogate among minimizers. For example $f(w_1,w_2)=(w_1-1)^2/2$ has $S=\{1\}\times\R$, on which $2(|w_1|+|w_2|)$ varies. A nearest point is a selected reference, not a canonical minimum-energy representative.

For a constrained or local version, replace $S$ explicitly by a nonempty closed set of minimizers of $f$ on a specified domain $K$, with common value $f_K$. If~\eqref{eq:qg} is assumed on $\mathcal U\cap K$ with $f_K$ and that $S$, the identical proof gives~\eqref{eq:stability} only for $\theta\in\mathcal U\cap K$. This is a separate hypothesis; it says nothing about infeasible points or other local minima.

\begin{corollary}[Strong convexity as a special case]\label{cor:sc}
Let $f$ be differentiable and $\mu$-strongly convex on a convex neighborhood $\mathcal U$ of a stationary point $\theta_*$, with $\nabla f(\theta_*)=0$. Then for $\theta\in\mathcal U$,~\eqref{eq:stability} holds with reference $\theta_*$, $\gamma=\mu$ and gap $f(\theta)-f(\theta_*)$.
\end{corollary}
\begin{proof}
The first-order strong-convexity inequality at $\theta_*$ gives
$f(\theta)-f(\theta_*)\ge(\mu/2)\norm{\theta-\theta_*}_2^2$. Apply Lemma~\ref{lem:perturbation}. For a constrained minimizer on a convex feasible domain, the same conclusion for feasible points follows if the variational inequality $\ip{\nabla f(\theta_*)}{\theta-\theta_*}\ge0$ holds there.
\end{proof}

\paragraph{Sharpness and limits.}
For $f(\theta)=(\gamma/2)\norm\theta_2^2$, $\theta_*=0$, zero bias perturbations, and each $W_\ell$ having $r_\ell$ nonzero singular values all equal to a common $c>0$, both sides of the first inequality in~\eqref{eq:stability} equal $2c\sum_\ell r_\ell$. This example verifies attainment of the constant in the stated general class. For a nontrivial equality case with $R(W,W^+)>0$, equality requires alignment in the reverse triangle inequality, common signs in the scalar summation, equality in the singular-value and layer Cauchy--Schwarz inequalities, no bias contribution to the distance, and equality in the growth bound at the point considered. Equal perturbation singular values alone are insufficient. This sharpness example is an elementary consequence, not a claim of new general stability theory.

A positive semidefinite quadratic $f(\theta)=\norm{B\theta-d}_2^2/2$ of rank-deficient $B$ has QG relative to its affine minimizer set with constant equal to the smallest positive eigenvalue of $B^\top B$ when $B\ne0$; it is not strongly convex on the full space. The projection is $\theta^+=\theta-B^\dagger(B\theta-d)$ and the gap equals $\norm{B(\theta-\theta^+)}_2^2/2$, proving the assertion by the singular values of $B$ (the least-squares residual is orthogonal to its range). In contrast, a selected positive eigenvalue of a neural-network Hessian establishes neither full-space strong convexity nor QG on a neighborhood. ReLU losses are also not automatically globally differentiable with Lipschitz gradient.

\section{Projection and first-order optimality}\label{sec:projection}
For $\tau\ge0$ define
\begin{equation}\label{eq:constraint}
 C_\tau=\{W:\Omega(W)\le\tau\},\qquad K_\tau=C_\tau\times\prod_\ell\R^{n_\ell}.
\end{equation}
The set $C_\tau$ is nonempty, compact and convex. Indeed, $\sum\norm{W_\ell}_F\le\sum\norm{W_\ell}_*\le\tau/2$, and the defining function is continuous and convex. The full parameter set $K_\tau$ is closed and convex but has unbounded bias directions. Consider
\begin{equation}\label{eq:problem}
 \min_{(W,b)\in K_\tau} f(W,b),
\end{equation}
where the chosen objective (data loss or an explicitly specified regularized loss) must be recorded. We do not assume that a minimum of~\eqref{eq:problem} exists merely because $W$ is budgeted.

\begin{proposition}[Joint Euclidean projection]\label{prop:projection}
Let $Y_\ell=Q_\ell\diag(\sigma_\ell)V_\ell^\top$ be reduced SVDs, including all $d_\ell$ singular values, and let $\sigma$ be their concatenation. The unique solution of
\begin{equation}\label{eq:projection}
 \min_{Z\in C_\tau}\frac12\sum_\ell\norm{Z_\ell-Y_\ell}_F^2
\end{equation}
is $Z_\ell=Q_\ell\diag((\sigma_\ell-t)_+)V_\ell^\top$ with the following conventions.
\begin{enumerate}
\item If $\Omega(Y)\le\tau$, choose $t=0$ and $Z=Y$.
\item If $0<\tau<\Omega(Y)$, $t\in(0,\max_i\sigma_i)$ is the unique solution of $h(t)=\sum_i(\sigma_i-t)_+=\tau/2$.
\item If $\tau=0$, then $Z=0$. Every $t\ge\max_i\sigma_i$ represents this same projection; the threshold is not unique. For $Y=0$, choose $t=0$.
\end{enumerate}
Repeated singular values and nonunique singular vectors do not affect the unique projected matrix.
\end{proposition}
\begin{proof}
For prescribed singular values $z_\ell\ge0$, the trace inequality $\ip{Z_\ell}{Y_\ell}\le\sum_i z_{\ell i}\sigma_{\ell i}$ implies
$\norm{Z_\ell-Y_\ell}_F^2\ge\sum_i(z_{\ell i}-\sigma_{\ell i})^2$, with equality when singular vectors are aligned. Thus~\eqref{eq:projection} reduces to
\[
 \min_{z\ge0,\;\boldsymbol1^\top z\le\tau/2}\tfrac12\norm{z-\sigma}_2^2.
\]
For an active positive radius, the KKT equations yield $z_i=(\sigma_i-t)_+$ and $\sum_i z_i=\tau/2$ with $t>0$. If the input is feasible the minimizer is itself. At radius zero the feasible set is a singleton. The function $h$ is continuous and strictly decreasing on $[0,\max\sigma]$ when $\max\sigma>0$, and identically zero on $[\max\sigma,\infty)$. The intermediate value theorem proves existence and uniqueness of the active positive-radius threshold. The original squared-distance objective is strictly convex on a convex set, proving uniqueness of $Z$ even with repeated singular values. The trace inequality is a standard matrix-analysis fact~\cite{Bhatia1997}; the scalar reduction is the familiar $\ell_1$-ball projection~\cite{Duchi2008L1proj}.
\end{proof}

Bisection evaluates $h$; it is not the only threshold solver. If $a_1\ge\cdots\ge a_D$ are the sorted concatenated singular values and $q=\tau/2>0$ in the active case, put
\[
 j_* =\max\left\{j:a_j>\frac{\sum_{i=1}^j a_i-q}{j}\right\},\qquad
 t=\frac{\sum_{i=1}^{j_*}a_i-q}{j_*}.
\]
This gives the same threshold~\cite{Duchi2008L1proj}. For objective $(\rho/2)\norm{Z-Y}_F^2$ and constraint $\Omega(Z)-\tau\le0$ with multiplier $\nu\ge0$, the singular-value stationarity equation is $\rho(z_i-\sigma_i)+2\nu=0$ when $z_i>0$. Therefore $t=2\nu/\rho$. If the constraint is written $\sum_i z_i\le\tau/2$, its multiplier is $2\nu$. The projection itself is independent of positive $\rho$; its multiplier scales with $\rho$.

\paragraph{Layer budgets and zero layers.}
For fixed radii $q_\ell\ge0$ with $\sum_\ell q_\ell=\tau/2$, the product set $\prod_\ell\{Z_\ell:\norm{Z_\ell}_*\le q_\ell\}$ is contained in $C_\tau$. Separate projections solve the product-set problem, generally with different thresholds. This restriction need not be worse for prediction, and no universal ordering of learning performance follows. With $Y_\ell\ne0$ and $q_\ell>0$, its individual nuclear-ball projection cannot be zero: it is either unchanged or has nuclear norm $q_\ell>0$. In a joint projection, however, an entire layer vanishes exactly when its largest input singular value is at most the common $t$. This is distinct from numerical SVD failure and from applying a fixed proximal threshold.

\paragraph{Approximate SVDs.}
The formula above requires all singular values relevant to the projection. It is not an exact projection of $Y$ merely because a truncated SVD is used. If $\widetilde Y$ satisfies $\norm{Y-\widetilde Y}_F\le\epsilon$ and $\widetilde Z=\Pi_{C_\tau}(\widetilde Y)$ is computed exactly, then nonexpansiveness gives
$\norm{\widetilde Z-\Pi_{C_\tau}(Y)}_F\le\epsilon$; $\widetilde Z$ remains feasible. Writing $d=\dist(Y,C_\tau)$, its projection-objective excess is at most $d\epsilon+\epsilon^2/2$, by the triangle inequality. For an arbitrary approximate output $\widehat Z$ within $\epsilon$ of the exact projection, feasibility is only bounded by $(\Omega(\widehat Z)-\tau)_+\le2\sqrt D\epsilon$. These estimates require a justified tail or output-error bound. Rounded arithmetic alone supplies no rigorous enclosure. The evaluation specified here uses full SVDs before studying approximations.

\begin{proposition}[Spectral-ball projection and scaling]\label{prop:spectral}
For $a\ge0$ and $Y=Q\diag(\sigma)V^\top$, the Frobenius projection onto $\{Z:\norm Z_2\le a\}$ is
\begin{equation}\label{eq:spectral}
 \Pi(Y)=Q\diag(\min(\sigma_i,a))V^\top.
\end{equation}
It generally differs from $\min(1,a/\sigma_{\max}(Y))Y$, interpreted as zero for $Y=0$ when needed. Positive uniform scaling of a nonzero matrix leaves its normalized singular-value distribution and entropy effective rank unchanged.
\end{proposition}
\begin{proof}
Use the same trace inequality as above to reduce to independently projecting each singular value onto $[0,a]$. For $Y=\diag(2,1)$, $a=1$, clipping gives $\diag(1,1)$ at squared distance $1$, while rescaling gives $\diag(1,1/2)$ at squared distance $5/4$. For any $c>0$, $c\sigma_i/\sum_j c\sigma_j=\sigma_i/\sum_j\sigma_j$, proving the last statement.
\end{proof}
Spectral normalization~\cite{Miyato2018SpectralNorm} parametrizes effective weights by division by an estimate of their largest singular value; it is not~\eqref{eq:spectral}. Measurements must use the weights actually applied in the forward pass. Power iteration can introduce additional approximation error. Spectral bounds $\norm{W_\ell}_2\le a_\ell$ do give the dimensional bound
\[
 \Omega(W)\le2\sum_\ell d_\ell a_\ell.
\]
Likewise $\norm{W_\ell}_F\le c_\ell$ implies $\Omega(W)\le2\sum_\ell\sqrt{d_\ell}c_\ell$. These bounds need not match a prescribed $\tau$. A penalty such as weight decay does not impose a fixed norm cap at every step, but with a lower-bounded data loss and a known regularized-objective sublevel it can also yield quantitative norm bounds. Blanket claims that SN or WD provides no control of $\Omega$ are therefore incorrect.

\begin{proposition}[Necessary optimality conditions]\label{prop:kkt}
Suppose $f$ is differentiable in a neighborhood of a local minimizer $\theta_*=(W_*,b_*)$ of~\eqref{eq:problem}. Then
\begin{equation}\label{eq:normal_kkt}
 0\in\nabla_W f(\theta_*)+N_{C_\tau}(W_*),\qquad \nabla_b f(\theta_*)=0,
\end{equation}
where $N_C(x)=\{v:\ip v{z-x}\le0\text{ for every }z\in C\}$ for $x\in C$. If $\tau>0$, there is $\nu_*\ge0$ such that
\begin{equation}\label{eq:scalar_kkt}
 \Omega(W_*)\le\tau,\quad \nu_*(\Omega(W_*)-\tau)=0,\quad
 0\in\nabla_{W_\ell}f(\theta_*)+2\nu_*\partial\norm{W_{\ell,*}}_*.
\end{equation}
For a compact positive-rank SVD $W_{\ell,*}=Q\Sigma V^\top$,
\begin{equation}\label{eq:subgradient}
 \partial\norm{W_{\ell,*}}_*=
 \{QV^\top+H:Q^\top H=0,\ HV=0,\ \norm H_2\le1\}.
\end{equation}
At zero this is the full spectral unit ball. These conditions are not sufficient for global optimality of a nonconvex network objective.
\end{proposition}
\begin{proof}
Every feasible line segment from a local minimizer has a nonnegative one-sided directional derivative, giving~\eqref{eq:normal_kkt}. For $\tau>0$, zero weights are strictly feasible for the continuous convex constraint; the normal-cone formula for this convex sublevel set is $N_{C_\tau}(W_*)=\{\nu G:\nu\ge0,G\in\partial\Omega(W_*),\nu(\Omega(W_*)-\tau)=0\}$, giving~\eqref{eq:scalar_kkt}. Formula~\eqref{eq:subgradient} is the nuclear-norm subdifferential~\cite{Watson1992Subdiff}: duality with the spectral norm characterizes a subgradient by $\norm G_2\le1$ and $\ip G W=\norm W_*$. Attaining this equality on the nonzero singular subspaces forces the $QV^\top$ component and leaves a spectral contraction on the orthogonal complements. For $W=0$ only $\norm G_2\le1$ remains. At $\tau=0$, use~\eqref{eq:normal_kkt} with $C_0=\{0\}$; the strict-feasibility argument is unavailable.
\end{proof}

\section{A split algorithm and finite-iteration certificates}\label{sec:algorithm}
Introduce $Z\in C_\tau$ and the equality $W=Z$. With scaled dual variable $u$ (never used for singular vectors), fixed $\rho>0$ and indicator $I_C$, the scaled augmented Lagrangian is
\begin{equation}\label{eq:augmented}
 \mathcal L_\rho(W,b,Z,u)=f(W,b)+I_{C_\tau}(Z)
 +\frac\rho2\norm{W-Z+u}_F^2-\frac\rho2\norm u_F^2.
\end{equation}
This is the equality-constrained splitting of~\eqref{eq:problem}. We do not use the expression $f+\nu g+(\rho/2)g^2$ as an augmented Lagrangian for the original inequality $g\le0$, because that expression penalizes strictly feasible points too.

\begin{algorithm}[htbp]
\caption{Exact ADMM step for a global nuclear-norm budget}\label{alg:admm}
\begin{algorithmic}[1]
\Require $\tau\ge0$, $\rho>0$, initial $(W^0,b^0)$; set $Z^0=\Pi_{C_\tau}(W^0)$, $u^0=0$.
\For{$k=0,1,\ldots$}
\State Choose an attained minimizer
\Statex\hspace{1em} $ (W^{k+1},b^{k+1})\in\argmin_{W,b}\bigl[f(W,b)+\frac\rho2\norm{W-Z^k+u^k}_F^2\bigr]$.
\State $Y_\ell\gets W_\ell^{k+1}+u_\ell^k$ for every layer.
\State $Z^{k+1}\gets\Pi_{C_\tau}(Y)$ using one threshold on all layer singular values (Proposition~\ref{prop:projection}).
\State $u^{k+1}\gets u^k+W^{k+1}-Z^{k+1}$.
\State Record $r^{k+1}=W^{k+1}-Z^{k+1}$ and $s^{k+1}=\rho(Z^{k+1}-Z^k)$.
\State Measure loss, predictions and $\Omega$ separately at $(W^{k+1},b^{k+1})$ and $(Z^{k+1},b^{k+1})$.
\EndFor
\end{algorithmic}
\end{algorithm}
The exact algorithm presupposes an attained solution of its first subproblem; without bias coercivity, even this requires care. A finite number of minibatch SGD steps, with or without momentum, is an inexact variant. It must record the inner update rule and projection frequency. Projected SGD, in contrast, applies $W^+=\Pi_{C_\tau}(W-\eta\widehat\nabla_W f)$ and updates biases separately. It is another way to obtain feasible postprojection iterates, so feasibility is not exclusive to ADMM.

If $\rho$ is changed after an iteration, first form $u_{\rm temp}=u^k+r^{k+1}$ with the old $\rho_k$, then set $u^{k+1}=(\rho_k/\rho_{k+1})u_{\rm temp}$ to preserve the unscaled multiplier $y=\rho u$. The actual ratio must be used even when $\rho$ is clipped at bounds. A changing $\tau_k$ changes the feasible problem: certificates must state whether they concern the current or final budget. Neither adaptation is part of the fixed-parameter exact algorithm above.

\begin{proposition}[Finite-iteration feasibility]\label{prop:feasibility}
For arbitrary $W,Z$ and $\tau\ge0$, with $v_\tau(X)=(\Omega(X)-\tau)_+$,
\begin{align}\label{eq:feasibility}
 v_\tau(W)&\le v_\tau(Z)+2\sum_\ell\norm{W_\ell-Z_\ell}_*\\
 &\le v_\tau(Z)+2\sqrt{R(W,Z)}\norm{W-Z}_F
 \le v_\tau(Z)+2\sqrt D\norm{W-Z}_F.\nonumber
\end{align}
In particular exact feasible $Z$ removes the first term.
\end{proposition}
\begin{proof}
Since $\Omega(W)-\tau\le\Omega(Z)-\tau+|\Omega(W)-\Omega(Z)|$, monotonicity and subadditivity of the positive part give the first bound via Lemma~\ref{lem:perturbation}. The remaining bounds follow from the same lemma.
\end{proof}
The direct nuclear-residual sum is the tightest of these computable upper bounds; $R$ exploits known residual ranks, whereas $D$ avoids rank estimation. A small Frobenius residual may still give a loose bound in a high-dimensional network. These are mathematical inequalities; their floating-point evaluations with stated tolerances are numerical certificates, not interval-validated proofs. A feasible $Z$ alone says nothing about exact feasibility of $W$ unless $r=0$.

\begin{proposition}[A posteriori stationarity at a feasible point]\label{prop:stationarity}
Consider one step with $\rho$ fixed during that step and exact $Z^+=\Pi_{C_\tau}(W^++u)$, $u^+=u+W^+-Z^+$. Put
\[
 x=(W^+,b^+),\quad z=(Z^+,b^+),\quad
 r=W^+-Z^+,\quad s=(\rho(Z^+-Z^{\rm old}),0),
\]
\[
 e=\bigl(\nabla_W f(x)+\rho(W^+-Z^{\rm old}+u),\ \nabla_b f(x)\bigr).
\]
Suppose $f$ is differentiable at $x,z$ and
$\norm{\nabla f(z)-\nabla f(x)}_2\le L\norm{x-z}_2$
(for example its gradient is $L$-Lipschitz on a region containing their segment). Then
\begin{equation}\label{eq:stationarity}
 \dist\bigl(0,\nabla f(z)+N_{K_\tau}(z)\bigr)
 \le\norm e_2+\norm s_2+L\norm r_F.
\end{equation}
\end{proposition}
\begin{proof}
Projection optimality implies $\rho(W^++u-Z^+)=\rho u^+\in N_{C_\tau}(Z^+)$. Thus $n=(\rho u^+,0)\in N_{K_\tau}(z)$, and direct substitution gives $\nabla f(x)+n=e-s$. Consequently
\[
 \dist(0,\nabla f(z)+N_{K_\tau}(z))
 \le\norm{\nabla f(z)+n}_2
 \le\norm e_2+\norm s_2+L\norm r_F.
\]
The cone is evaluated at the feasible point $z$, not at $x$.
\end{proof}
An exact differentiable first-subproblem minimizer gives $e=0$, but practical SGD does not. A minibatch gradient cannot be silently substituted for the full gradient in $e$. If a certified bound $\norm{\widehat g-\nabla f(x)}_2\le\delta$ is available, the same argument using $\widehat g$ in $e$ adds $\delta$ to the right side. For a feasible inexact projection $\widehat Z$, if
$\dist(\rho(W^++u-\widehat Z),N_{C_\tau}(\widehat Z))\le\eta$, the same proof adds $\eta$ to~\eqref{eq:stationarity}, using the corresponding $r,s,e,z$. Feasibility error alone does not bound this normal-cone defect. If the approximate output is infeasible, this statement does not apply there. For ReLU networks, a justified smooth region and $L$ must be supplied; otherwise report the ingredients as diagnostics and do not label their sum a validated smooth stationarity certificate.

\paragraph{Why block descent is insufficient for convergence.}
At fixed $W,b,Z,\rho$, write $r=W-Z$. Expanding~\eqref{eq:augmented} gives a dual-dependent term $\rho\ip u r$. Therefore
\begin{equation}\label{eq:dual_increase}
 \mathcal L_\rho(W,b,Z,u+r)-\mathcal L_\rho(W,b,Z,u)=\rho\norm r_F^2.
\end{equation}
The dual step increases this Lagrangian. Primal block descent alone cannot imply a decreasing telescopic bound after the dual update. Exact projection, smoothness and bounded iterates therefore do not constitute the earlier claimed proof of residual convergence. We do not claim monotone loss, monotone residuals, a nonconvex ADMM convergence theorem, or a linear convergence rate for the implemented stochastic algorithm.

For comparison, the standard convex ADMM result in Boyd et al.\ (Sections 3.2--3.3)~\cite{Boyd2011ADMM} uses closed proper convex objectives, existence of a saddle point, exact subproblem solutions and fixed $\rho$. These assumptions are not established for the neural-network implementation. The controlled quadratic experiments can study that regime, while~\eqref{eq:feasibility} and~\eqref{eq:stationarity} remain finite-iteration statements. Small measured residuals in a finite run do not prove an asymptotic limit.

\section{Penalty bounds and claims that require extra hypotheses}\label{sec:penalty}
Define the separate hinge-penalty objective
\[
 F_\lambda(\theta)=f(\theta)+\lambda v_\tau(W),\qquad\lambda>0.
\]
It is generally nonsmooth and is not the objective of the ADMM first step. A nuclear penalty $f+\lambda\Omega$ is yet another objective, without the hinge. Smooth-gradient PL theory cannot be applied to any of these by changing names. Exact penalization is classical~\cite{HanMangasarian1979}; the following elementary bounds state precisely what is used here.

\begin{proposition}[Valid penalty-violation bounds]\label{prop:penalty}
If $a$ is a genuine lower bound for $f$ on its domain, then
\begin{equation}\label{eq:penalty_lower}
 \lambda v_\tau(W)\le F_\lambda(\theta)-a.
\end{equation}
Alternatively, assume there are a globally constrained optimal point $\theta_*$, value $f_*=f(\theta_*)$ and $\nu_*\ge0$ with feasibility, complementarity and the global saddle inequality
\begin{equation}\label{eq:saddle}
 f(\theta)+\nu_*(\Omega(W)-\tau)\ge f_*\quad\hbox{for every }\theta.
\end{equation}
If $\lambda>\nu_*$, then $\min F_\lambda=f_*$ and
\begin{equation}\label{eq:penalty_saddle}
 (\lambda-\nu_*)v_\tau(W)\le F_\lambda(\theta)-f_*.
\end{equation}
Every global minimizer of $F_\lambda$ is then a constrained global minimizer.
\end{proposition}
\begin{proof}
The first inequality is immediate from $f\ge a$. For the second, if $g=\Omega(W)-\tau\ge0$, then~\eqref{eq:saddle} gives $F_\lambda-f_*\ge(\lambda-\nu_*)g$. If $g<0$, it gives $F_\lambda=f\ge f_*-\nu_*g\ge f_*$. Equality $F_\lambda(\theta_*)=f_*$ follows from feasibility. Thus the minimum is $f_*$, and equality at any minimizer forces $v_\tau=0$ when $\lambda>\nu_*$. Its objective value is then $f_*$.
\end{proof}
Global saddle conditions follow from convexity and suitable constraint qualification, but are not consequences of local neural-network KKT conditions. Neither bound warrants a violation estimate using $\lambda^{-1}(F_\lambda-F_\lambda^*)$ in general. For
\[
 f(w)=\tfrac12(w-2)^2,\quad\Omega(w)=2|w|,\quad\tau=2,\quad\lambda=1,
\]
$F_\lambda$ has minimizer $w_*=1$ and minimum $1/2$, but at $w=3/2$ its gap is $5/8$ while the violation is $1$. Even a feasible penalized optimum does not repair the claimed bound; the data loss can decrease outside the constraint. For $0<\lambda<1/2$ in this example, the minimizer $2-2\lambda>1$ is itself infeasible.

If~\eqref{eq:saddle} holds, $\Omega(W_*)=\tau$ and $\nu_*>0$, then for feasible $\theta$,
\[
 0\le\Omega(W_*)-\Omega(W)\le\frac{f(\theta)-f_*}{\nu_*}.
\]
This is a valid replacement for a multiplier-based gap statement. An active constraint does not imply a positive multiplier: minimizing $(w-1)^2/2$ over $|w|\le1$ has active solution $1$ and multiplier zero. Without convexity or a global saddle property, the displayed bound need not follow from KKT: $f(w)=1-w^3$, $\Omega=2|w|$, $\tau=2$, $w_*=1$, $\nu_*=3/2$ satisfies KKT, yet $w=0$ would require $2\le2/3$.

Finally, a recursion $\Delta_{k+1}\le q\Delta_k+\epsilon_k$, $0<q<1$, yields
\[
 \Delta_k\le q^k\Delta_0+\sum_{j=0}^{k-1}q^{k-1-j}\epsilon_j.
\]
Summable errors alone do not imply a geometric rate: for equality and $\epsilon_j=(j+1)^{-2}$, $\Delta_k\ge k^{-2}$. Even for an unrelated smooth PL objective, the exact gradient-descent decrease with step $\alpha\in(0,2/L)$ gives the factor $1-2\alpha\mu(1-\alpha L/2)$, not the unrestricted factor $1-\alpha\mu$. Since those hypotheses and that iteration are not established here, no PL rate is retained for the neural-network algorithm.

\section{Verification, historical evidence and prospective evaluation}\label{sec:evidence}
\subsection{Small mathematical checks completed in this revision}
The accompanying script \texttt{comprobaciones/verificar\_matematica.py} executes only small CPU calculations in double precision. It does not load the training notebook or train a neural network. With NumPy 1.26.4, Python 3.11.8 and random seed 20261013, all 227 recorded checks passed. The complete machine-readable output and actual execution log accompany the revision. These checks exercise the formulas and counterexamples; the proofs above establish the mathematical claims.

The checks include the path example $E=2\sqrt2$, $\Omega=4$; 100 small signed layered graphs satisfying~\eqref{eq:bilateral} within floating-point tolerances; joint-projection feasibility and idempotence for zero, positive, interior and repeated-spectrum cases; comparison of sorting, bisection and an independent constrained quadratic solver on the singular values; and spectral clipping versus rescaling. They also check the QG equality example and an explicit rank-deficient least-squares projection, 50 finite-feasibility instances, the full-space stationarity identity including bias gradients, the positive dual Lagrangian increment, both penalty and multiplier counterexamples, polynomial error accumulation, and ReLU rescaling. These are numerical checks with tolerances, not rigorous rounding-error enclosures or evidence of predictive superiority.

\subsection{Status of the historical neural-network experiments}
The previous manuscript reports MNIST experiments on three ReLU MLPs, listed in Table~\ref{tab:architectures}, and compares standard training, weight decay, nuclear penalization, spectral normalization and an ADMM variant. The original sources, PDF, figures and CSV files are preserved unchanged. Their outcomes are historical exploratory observations. Their previous interpretations as convergence proofs, validation of strong-convexity hypotheses, and exclusive ADMM feasibility guarantees are withdrawn.

\begin{table}[htbp]
\centering
\caption{Architectures retained for the prospective evaluation. Biases are trainable. Width changes accompany depth changes, so these models do not isolate depth.}
\label{tab:architectures}
\begin{tabular}{llr}
\toprule
Name & Layer widths & Parameters \\
\midrule
3L & 784--256--128--10 & 235,146\\
6L & 784--256--128--64--64--32--10 & 248,682\\
9L & 784--256--128--128--64--64--32--32--16--10 & 266,618\\
\bottomrule
\end{tabular}
\end{table}

A static inspection of the public experimental repository at commit \texttt{186e27400f94} (full hash recorded in the audit) identified measurement and provenance issues. The spectral-normalization diagnostics iterate over named parameters, which can refer to original unconstrained parameters rather than effective forward-pass matrices; a separate spectral reconstruction also mismatches parametrized names. The historical curvature calculation selects the smallest positive Hessian eigenvalue, clips negative loss gaps, and uses thresholded perturbation ranks. These operations do not verify the hypotheses or constants of Theorem~\ref{thm:qg}. The projection utilities include fallbacks that replace failed SVDs by zeros, while warnings are globally suppressed. Budget satisfaction is counted using a one-percent allowance in parts of the notebook, rather than exact feasibility or the explicit tolerances proposed below. The quantity called validation uses samples from the official test partition. These issues prevent the existing summaries from serving as confirmatory evidence for the revised claims.

Configurations are repeated across notebook experiments without run hashes establishing independence; the earlier count of twenty independent ADMM runs is unsupported. The original discussion mixes three-seed and five-seed summaries. The ADMM penalty adaptation rescales the dual by an assumed factor even when the penalty is clipped, rather than always using the actual ratio. In contrast, the inspected nuclear-penalty path uses differentiable tensor SVD values; the specific suspicion that this penalty necessarily loses its gradient through a CPU transfer is not supported by that code path. The evidence report records cell and line locations and separates observed defects from risks awaiting execution-level checks.

Table~\ref{tab:historical} summarizes which historical material can be reused and which interpretations or measurements require replacement.
\begin{table}[htbp]
\caption{Disposition of historical evidence. Figure numbers refer to the old manuscript, not to new results.}
\label{tab:historical}
\small
\begin{tabularx}{\linewidth}{>{\raggedright\arraybackslash}p{0.19\linewidth}>{\raggedright\arraybackslash}X >{\raggedright\arraybackslash}X}
\toprule
Historical material & Status & Required action \\
\midrule
Architectures and method definitions & Reusable after correcting conventions & Preserve shapes, seeds and provenance; specify objectives and effective weights.\\
Figure 1 and curvature-based ratios & Invalid as validation of SC/QG & Replace with known minimizers and growth constants. A ratio above one is a failed check, not a vacuous bound.\\
Figure 2 and residual summaries & Pending numerical audit; convergence interpretation invalid & Identify $W$ versus $Z$, budgets, tolerances and distinct runs. The plotted primal residual is not monotone.\\
Figures 3--4 and accuracy tables & Pending measurement and split audit & Use independent validation/test splits, paired budgets and all-method feasibility reporting.\\
Figure 5, single-seed layer allocation & Exploratory only; claimed annihilation mechanism invalid & Repeat across seeds; log spectra and gradients. A positive-radius single-layer projection does not annihilate a nonzero matrix.\\
Figures 6--7 and SN spectral tables & SN measurement/interpretation invalid; other entries pending audit & Measure effective weights. Uniform scaling does not flatten normalized singular values.\\
Timing and compression claims & Pending controlled measurement & Separate hardware effects and calibration costs; demonstrate stored-model compression if claimed.\\
\bottomrule
\end{tabularx}
\end{table}
No corrected neural-network accuracy or timing numbers are substituted here. Reusing a historical number requires a traceable configuration, dataset split, checkpoint, matching effective weights and an independently recomputed metric. Preserving original files does not endorse their conclusions.

\subsection{Prospective evaluation tied to the mathematical statements}
The complete implementation specification, experimental route and claim-to-evidence matrix accompany this revision. They define seven stages, E0--E6. E0 must pass before training: test the full joint projection, an independent small optimization solver, spectral projection, factor-of-two conventions, effective SN weights, nuclear-penalty gradients and non-silent failure handling. E1 uses a strongly convex matrix quadratic and rank-deficient least squares with explicit optimal sets, projections and QG constants. It does not estimate a neural-network growth constant from a selected Hessian eigenvalue.

E2 compares STD, WD, a nuclear penalty, joint ADMM and projected SGD with momentum on MNIST~\cite{lecun1998gradient} and Fashion-MNIST~\cite{Xiao2017FashionMNIST}, using Table~\ref{tab:architectures} and five predefined seeds $\{42,123,456,789,1024\}$. SN is retained only with correct effective-weight measurements. For each dataset, architecture and seed, one frozen standard-training reference defines $\tau=\alpha\Omega(W_{\rm STD,ref})$ with $\alpha\in\{0.6,0.8,1.0\}$. The two constrained methods alone require $2\times3\times3\times5\times2=180$ runs. References, WD/NUC calibration, SN and ablations add separate costs, deduplicated by configuration hashes. Hyperparameters use training and validation data, with the official test set reserved for final evaluation. A pilot measures cost before the principal run allocation is approved; no runtime or GPU configuration is assumed here.

E3 records $\Omega(W)$, $\Omega(Z)$, current and final violations, primal and dual residuals, $\rho$ and projection frequency. The primary ADMM model is the actual $W$ iterate; $(Z,b)$ and $(\Pi_{C_\tau}(W),b)$ are separately labelled secondary models. Accuracy and spectral measurements always refer to the same weights. Exact projection makes $Z$ feasible, while Proposition~\ref{prop:feasibility} bounds possible violation of $W$. The smooth stationarity certificate is assessed only when its full-gradient and smoothness assumptions are justified. Warmup, budget transitions and fixed-budget phases are reported separately.

E4 compares global projection with per-layer budgets proportional to parameter counts or reference nuclear norms, with an optional uniform allocation. The feasible sets differ. Several seeds, layer spectra, gradients and a longer-horizon or width-controlled diagnostic are required before attributing failures to depth or allocation. E5 checks $E(A)/\Omega$ on small graphs and selected checkpoints, treating $\Omega=0$ separately, and verifies ReLU rescaling without interpreting it as preservation of a weight-decayed objective. E6 measures total and component time on the same Linux/CUDA system, with synchronization and warmup. Actual storage and predictive effects of a factorized or truncated model are needed for any compression claim; effective rank alone is insufficient.

Numerical feasibility uses a preregistered absolute-plus-relative tolerance, with raw violation retained. For E0 float64 checks the feasibility tolerance is $10^{-10}+10^{-10}\max(1,\tau)$. For the primary float32 training outputs, spectral measurements are reevaluated in float64 and feasibility uses $10^{-6}+10^{-5}\max(1,|\tau|)$. Any amendment must be justified through E0 before test evaluation. No NaN, failed SVD or failed training run is converted silently to a successful zero-energy result. Results include seed-level dispersion, paired differences, failures and retry histories. Standard deviation is not a confidence interval, and five seeds allow limited precision. The replicate is a run/seed, not an epoch or image.

\section{Discussion and conclusions}\label{sec:conclusions}
The revised analysis distinguishes the exact bipartite block identity, a two-sided comparison for the full consecutive-layer graph, and a local loss-gap bound for the additive surrogate under quadratic growth. The stability argument is a direct consequence of matrix-norm inequalities and the stated growth hypothesis. Its assumptions and direction are essential: a small spectral budget does not itself imply low loss, and equivalent ReLU parametrizations may have different budgets.

A global nuclear-norm budget admits an exact common-threshold Euclidean projection. Separate layer budgets, spectral clipping, spectral normalization and nuclear penalization are distinct operations, each potentially useful for an appropriately specified objective. Projection guarantees feasibility of the point projected; it does not by itself prove convergence of a stochastic split algorithm or guarantee predictive quality.

The finite-iteration residual certificate and conditional stationarity estimate provide diagnostics independent of an asymptotic convergence claim. The revised work establishes these statements in the main text and records small numerical checks. It does not establish convergence of the neural-network ADMM implementation, intrinsic functional complexity control, electrical-energy reduction or deployment compression.

The essential neural-network evaluation remains pending. If ADMM outperforms matched-budget projected SGD with momentum and calibrated penalties under the specified protocol, that would support a bounded empirical advantage for those settings. Similar outcomes would support feasibility with comparable predictive performance. Worse accuracy, greater cost, depth-specific failures or feasibility confined to $Z$ must be reported with equal prominence. None changes the algebraic results, but each changes the admissible applied conclusion. A submission decision requires completed experimental evidence and author agreement on the contribution and declarations.

\section*{Data and code availability}
The historical repository is \url{https://github.com/cero1979/Ensayos-Energy-Graph}. The revision records a static snapshot and does not claim that its notebook currently reproduces every published table. The accompanying package contains mathematical checks, an execution log, audit, implementation specification and prospective protocol. A versioned archive of corrected experimental code, configurations, raw metrics, data-split hashes and checkpoints remains to be prepared by the experimental collaborator before submission. No new training results were generated for this revision.

\section*{Author declarations requiring confirmation}
The original manuscript lists Carlos Ram\'irez Ovalle, Abel Alvarez and Andr\'es J. Castrill\'on V\'asquez at Pontificia Universidad Javeriana Cali. Names, affiliations, contributions, funding, acknowledgments and competing interests require confirmation by all authors. The full postal address of the affiliation must be supplied and verified for the title page before submission. The former draft's broad equal-contribution statement has not been treated as newly verified. No editorial submission or author approval is implied by this revision.

\section*{Declaration of generative AI and AI-assisted technologies in the manuscript preparation process}
\textit{Proposed disclosure, pending author review and approval.} During the preparation of this revision, the authors used OpenAI Codex to assist with mathematical auditing, drafting and restructuring proofs, checking small numerical examples, reviewing code statically, preparing an experimental specification, verifying bibliographic information and revising English text. The authors must review and verify the mathematical arguments, code-related observations, references and final text before submission, edit the material as needed, and take full responsibility for the content. This proposed statement describes substantive assistance and does not limit it to language correction.

\bibliographystyle{elsarticle-num}
\bibliography{reference}
\end{document}
